\subsubsection{服务窗口与两类中继资源的连接}
几何覆盖和时间覆盖分别承担不同约束。候选位置属于$\mathcal A_h$，表示能够满足需求区间$h$的接入和回传；实际指派$z_{hq}=1$后，还要求其在线窗口包含运输的绝对需求区间。几何条件成立而中继未到位时，不能向运输任务提供服务。

中继任务同时占用实体机和能源组件。若$q'$使用执行$q$的同一架中继机，需要等到$q$返航并完成周转；若使用同一能源组件，则需要等到该组件充满。两种接续分别为
\begin{equation}
 s_{q'}^R\ge e_q^R+t^{turn},\qquad
 s_{q'}^R\ge e_q^R+c_q^R,
 \label{eq:q3_two_relay_resources}
\end{equation}
各式仅在对应资源相同且$q$为前项时启用。题目允许同型资源复用，因此同一中继机可更换组件，已恢复组件也可供另一实体机使用。

服务结束$b_q$同时决定返航时刻、悬停能耗和能源恢复。延长窗口可能覆盖更多运输需求，却推迟后续资源；删去不必要的早到悬停，则可能同时节能并改善接续。将服务起止保留为连续变量，才能把通信保障真正纳入调度，而不是在运输时间表外另放若干静态点。

\subsubsection{分层计算中的决策接口}
候选重定位负责新的位置、高度和传播关系，联合整数模型负责任务选择、提供方和资源路径，固定结构线性规划负责最后的连续时刻。空间计算不直接塞进整数规划的代数约束，而是生成可审查的可服务关系；给定离散选择后，时间条件尽量交给线性优化精确处理。

这种接口还说明了扰动后的处理层次：只改变中继最早上线时刻，可复用几何并修复时序；传播门限改变，需要重检距离根与可服务区间；候选不再能覆盖关键需求，则回到提供方或位置选择层。每次调整都以同一运输需求和物理规则为基础，避免不同模块分别“优化”后得到互不兼容的日程。
